\documentclass[../DoAn.tex]{subfiles}
\begin{document}

\section{Tổng quan giải pháp}
\label{section:3.1}

\subsection{Định hướng tiếp cận tổng thể}
\label{subsection:3.1.1}
\begin{itemize}[nosep, leftmargin=*]
    \item Microsoft GraphRAG làm nền tảng, tùy biến cho domain luật lao động VN
    \item Ba trụ cột: (1) Knowledge Graph hai lớp, (2) Dual search + multi-hop + temporal filter, (3) Ứng dụng chat + phân tích HĐLĐ
    \item Baselines thí nghiệm: Basic/Naive RAG, Local Search, Global Search, Local + Multi-hop
\end{itemize}

\subsection{Phạm vi corpus và giả định thiết kế}
\label{subsection:3.1.2}
\begin{itemize}[nosep, leftmargin=*]
    \item 7 văn bản, 581 Điều --- nguồn \texttt{data/txt/}
    \item Giả định: văn bản đã chuẩn hóa \texttt{.txt}; LLM extraction có thể sai $\rightarrow$ bù bằng Lớp 1 deterministic
    \item Phạm vi ngoài: bản án, luật dân sự/hình sự (chỉ test mẫu, chưa index)
\end{itemize}

\section{Xây dựng đồ thị tri thức cho văn bản luật lao động}
\label{section:3.2}

\subsection{Mục tiêu xây dựng đồ thị tri thức}
\label{subsection:3.2.1}
\begin{itemize}[nosep, leftmargin=*]
    \item Chuyển văn bản phi cấu trúc $\rightarrow$ tri thức có cấu trúc phục vụ truy vấn và suy luận
    \item Mọi câu trả lời truy vết được về Điều/Khoản cụ thể
\end{itemize}

\subsection{Thiết kế ontology hai lớp}
\label{subsection:3.2.2}
\begin{itemize}[nosep, leftmargin=*]
    \item \textbf{Lớp 1 (deterministic):} \texttt{VanBan}, \texttt{Chuong}, \texttt{Dieu}, \texttt{Khoan}, \texttt{Diem}; quan hệ \texttt{contains}, \texttt{cites}, \texttt{amends}, \ldots
    \item \textbf{Lớp 2 (LLM):} \texttt{ChuThe}, \texttt{HanhVi}, \texttt{CheTai}, \texttt{TienLuong}, \ldots; quan hệ \texttt{entitles}, \texttt{obligates}, \texttt{penalizes}, \ldots
    \item \textbf{Lớp 1.5:} Canonicalization alias ``Điều 35'' $\rightarrow$ \texttt{BLLĐ\_2019\_Điều\_35}
\end{itemize}

\subsection{Quy trình chuẩn bị dữ liệu}
\label{subsection:3.2.3}
\begin{itemize}[nosep, leftmargin=*]
    \item Pipeline \texttt{scripts/01\_prepare\_data.py}: parse Phần/Chương/Mục/Điều/Khoản/Điểm
    \item Gán \texttt{norm\_type}, metadata hiệu lực
    \item Xuất JSONL + \texttt{metadata.json} + \texttt{normalized/}
    \item Thống kê: 581 Điều, 2069 Khoản, 1705 Điểm
\end{itemize}

\subsection{Chiến lược chunking}
\label{subsection:3.2.4}
\begin{itemize}[nosep, leftmargin=*]
    \item \texttt{chunking.size: 4000}, \texttt{overlap: 0} --- mục tiêu 1 Điều = 1 chunk
    \item Khoản là đơn vị pháp lý không được cắt đôi
    \item \texttt{input.text\_column: noi\_dung}; metadata giữ trong \texttt{raw\_data}
\end{itemize}

\subsection{Trích xuất entity và quan hệ (Lớp 2)}
\label{subsection:3.2.5}
\begin{itemize}[nosep, leftmargin=*]
    \item Prompt: \texttt{prompts/extract\_graph\_simple.txt}
    \item Quy tắc domain: phân biệt \texttt{HanhVi}/\texttt{CheTai}/\texttt{XuLyKyLuat}
    \item \texttt{max\_gleanings: 1} --- cân bằng chi phí/độ phủ
    \item Không extract Lớp 1 bằng LLM --- tránh duplicate node
\end{itemize}

\subsection{Merge graph và community detection}
\label{subsection:3.2.6}
\begin{itemize}[nosep, leftmargin=*]
    \item \texttt{02\_merge\_structural\_graph.py}: gộp L1 + output GraphRAG
    \item Community detection (Leiden) + \texttt{community\_report\_labor.txt}
    \item \texttt{summarize\_descriptions.max\_length: 500}
\end{itemize}

\section{Chiến lược truy vấn và hỏi đáp}
\label{section:3.3}

\subsection{Tổng quan truy xuất trong GraphRAG}
\label{subsection:3.3.1}
\begin{itemize}[nosep, leftmargin=*]
    \item Phân biệt vector retrieval vs graph-augmented retrieval
    \item Lựa chọn mode theo loại câu hỏi (single-hop vs tổng hợp)
\end{itemize}

\subsection{Local Search}
\label{subsection:3.3.2}
\begin{itemize}[nosep, leftmargin=*]
    \item Entity extraction từ query $\rightarrow$ retrieve text units + relationships
    \item Triển khai: \texttt{query/local\_search.py} $\rightarrow$ \texttt{ask\_local()}
    \item API: \texttt{POST /api/chat} với \texttt{mode: "local"}
    \item Phù hợp: tra cứu một Điều (LD001--LD004, LD013)
\end{itemize}

\subsection{Global Search}
\label{subsection:3.3.3}
\begin{itemize}[nosep, leftmargin=*]
    \item Map-reduce trên community reports
    \item Triển khai: \texttt{query/global\_search.py} $\rightarrow$ \texttt{ask\_global()}
    \item Phù hợp: câu hỏi tổng hợp (quyền NLĐ, loại hình HĐLĐ)
    \item Trade-off: latency cao hơn Local
\end{itemize}

\subsection{Basic Search --- baseline Naive RAG}
\label{subsection:3.3.4}
\begin{itemize}[nosep, leftmargin=*]
    \item Vector-only trên text units, không khai thác đồ thị
    \item Vai trò baseline trong thí nghiệm Chương 4
\end{itemize}

\subsection{Multi-hop Reasoning Engine}
\label{subsection:3.3.5}
\begin{itemize}[nosep, leftmargin=*]
    \item \texttt{query/multihop\_reasoning.py} --- class \texttt{VNLegalReasoningEngine}
    \item Chain templates: \texttt{violation}, \texttt{entitlement}, \texttt{procedure}
    \item Ví dụ: ``không trả lương'' $\rightarrow$ \texttt{obligates} $\rightarrow$ \texttt{penalizes} $\rightarrow$ NĐ 12/2022
    \item Tích hợp: phân tích HĐLĐ (\texttt{mapper\_optimized.py}, Mụ~\ref{subsection:3.5.2.1}); sắp mở rộng cho chat API
\end{itemize}

\subsection{Temporal Filter --- lọc hiệu lực văn bản}
\label{subsection:3.3.6}
\begin{itemize}[nosep, leftmargin=*]
    \item \texttt{query/temporal\_filter.py} --- \texttt{filter\_citations\_by\_effectiveness()}
    \item Metadata: \texttt{con\_hieu\_luc}, ngày ban hành, \texttt{huong\_dan\_boi}
    \item Tích hợp API: field \texttt{temporal\_warnings} trong \texttt{ChatResponse}
\end{itemize}

\subsection{Rule Validator}
\label{subsection:3.3.7}
\begin{itemize}[nosep, leftmargin=*]
    \item \texttt{query/rule\_validator.py} --- xác thực câu trả lời với quy tắc cứng
    \item Kiểm tra Điều bắt buộc, xác thực số liệu (lương tối thiểu, hạn thử việc)
    \item Sắp integrate vào pipeline query
\end{itemize}

\section{Tích hợp LLM và sinh câu trả lời}
\label{section:3.4}

\subsection{Vai trò LLM trong pipeline}
\label{subsection:3.4.1}
\begin{itemize}[nosep, leftmargin=*]
    \item Indexing: extraction, summarize, community report ($\sim$4000--5000 LLM calls)
    \item Querying: sinh câu trả lời từ ngữ cảnh retrieved
    \item LLM không phải nguồn tri thức --- chỉ tổng hợp và diễn đạt
\end{itemize}

\subsection{Thiết kế ngữ cảnh đầu vào}
\label{subsection:3.4.2}
\begin{itemize}[nosep, leftmargin=*]
    \item Context: text units + entity descriptions + community summaries
    \item Giới hạn context window; ưu tiên relevance
    \item Yêu cầu trích dẫn Điều/Khoản trong system prompt
\end{itemize}

\subsection{Chiến lược prompt và kiểm soát hallucination}
\label{subsection:3.4.3}
\begin{itemize}[nosep, leftmargin=*]
    \item Prompt trả lời: chỉ dựa trên ngữ cảnh; từ chối nếu thiếu thông tin
    \item Prompt extraction: quy tắc JSON output, entity type cố định
    \item Cache LLM (\texttt{data/labor-law/cache/}) --- tiết kiệm chi phí re-index
\end{itemize}

\section{Ứng dụng phân tích hợp đồng lao động}
\label{section:3.5}

\subsection{Mục tiêu và phạm vi}
\label{subsection:3.5.1}
\begin{itemize}[nosep, leftmargin=*]
    \item Kiểm tra điều khoản HĐLĐ so với BLLĐ 2019 + NĐ hướng dẫn
    \item Module: \texttt{unified-search-app/app/contract\_analysis/}
    \item Kiến trúc \textbf{hybrid}: rule-based (deterministic) + đồ thị \texttt{labor-law} (grounding) + LLM (phân tích khoản phức tạp)
    \item Phạm vi: HĐLĐ dạng text/PDF; không thay thế tư vấn pháp lý chính thức
\end{itemize}

\subsection{Pipeline phân tích HĐLĐ}
\label{subsection:3.5.2}
\begin{itemize}[nosep, leftmargin=*]
    \item \textbf{Bước 1 --- Tiếp nhận tài liệu}:
    \begin{itemize}[nosep, leftmargin=*]
        \item Đầu vào: PDF (text layer hoặc PaddleOCR khi scan), DOCX, TXT
        \item Đầu ra: \texttt{ContractDocument} gồm văn bản thô và metadata tệp
    \end{itemize}
    \item \textbf{Bước 2 --- Tách khoản}:
    \begin{itemize}[nosep, leftmargin=*]
        \item LLM phân đoạn HĐLĐ thành danh sách \texttt{Clause} (tiêu đề, tóm tắt, category, văn bản gốc)
        \item Phát hiện điều khoản bắt buộc còn thiếu (\texttt{compute\_missing\_mandatory})
    \end{itemize}
    \item \textbf{Bước 3 --- Kiểm tra quy tắc VR001--VR016}:
    \begin{itemize}[nosep, leftmargin=*]
        \item Regex + ngưỡng số học trên từng khoản; thời gian $<$ 5\,ms/khoản
        \item Căn cứ pháp lý gắn sẵn trong rule (ví dụ VR001: BLLĐ Điều~90 + NĐ~74/2024)
        \item Tham số \texttt{RuleContext}: vùng lương I--IV, \texttt{max\_probation\_days}
        \item Đây là lớp phát hiện vi phạm \textbf{ưu tiên} cho các trường hợp rõ ràng (lương, thử việc, giờ làm, phạt tiền NLĐ, \ldots)
    \end{itemize}
    \item \textbf{Bước 4 --- Phân loại khoản cần map đồ thị}:
    \begin{itemize}[nosep, leftmargin=*]
        \item Chỉ khoản thuộc nhóm phức tạp hoặc \texttt{UNKNOWN} mới gọi mapper
        \item Nhóm phức tạp: \texttt{SALARY}, \texttt{WORKING\_HOURS}, \texttt{PROBATION}, \texttt{TERMINATION}, \texttt{SOCIAL\_INSURANCE}, \texttt{PENALTY\_CLAUSE}, \texttt{UNILATERAL\_TERMS}, \texttt{WAIVER\_CLAUSE}, \texttt{NON\_COMPETE}, \texttt{DISPUTE\_RESOLUTION}
        \item Khoản đơn giản đã có kết quả rule đủ --- bỏ qua bước map để giảm latency
    \end{itemize}
    \item \textbf{Bước 5 --- Đối chiếu đồ thị pháp luật} --- xem Mục~\ref{subsection:3.5.2.1}
    \item \textbf{Bước 6 --- Phân tích LLM bổ sung}:
    \begin{itemize}[nosep, leftmargin=*]
        \item Batch review trên khoản phức tạp/chưa rõ; prompt nhận thêm \texttt{MappedLawSnippet} và tóm tắt rule issues
        \item \textbf{Bỏ qua LLM} nếu VR001--VR007 đã báo \texttt{VIOLATION} rõ --- tránh trùng lặp và tiết kiệm token
        \item LLM \emph{không} là nguồn tri thức pháp luật; chỉ đối chiếu khoản HĐ với ngữ cảnh đã retrieve
    \end{itemize}
    \item \textbf{Bước 7 --- Sinh báo cáo}:
    \begin{itemize}[nosep, leftmargin=*]
        \item Gộp rule issues + LLM issues; mức độ: \texttt{VIOLATION}, \texttt{HIGH\_RISK}, \texttt{MEDIUM\_RISK}, \texttt{COMPLIANT}
        \item Mỗi cảnh báo kèm căn cứ Điều/khoản khi có trong ngữ cảnh pháp lý
    \end{itemize}
    \item Tuỳ chọn: lưu phiên phân tích vào Neo4j (\texttt{neo4j\_store.py}, label \texttt{ContractSession}, TTL 24h) --- tách biệt đồ thị HĐ và đồ thị \texttt{labor-law}
\end{itemize}

\subsection{Tích hợp đồ thị \texttt{labor-law} trong phân tích HĐLĐ}
\label{subsection:3.5.2.1}
\begin{itemize}[nosep, leftmargin=*]
    \item \textbf{Vai trò}: đồ thị \emph{bổ trợ grounding} cho LLM --- \textbf{không thay thế} VR001--VR016
    \item \textbf{Nguồn dữ liệu}: artifacts GraphRAG tại \texttt{data/labor-law/} --- \texttt{entities.parquet}, \texttt{relationships.parquet}, \texttt{merged\_entities/relationships} (L1 + L2, xem Mục~\ref{subsection:3.2.6})
    \item \textbf{Mapper mặc định}:
    \begin{itemize}[nosep, leftmargin=*]
        \item \texttt{GraphLoader(root\_dir="data/labor-law")} load parquet đã merge
        \item \texttt{VNLegalReasoningEngine} (Mục~\ref{subsection:3.3.5}) --- chọn chain theo category khoản:
        \texttt{violation} (phạt, chấm dứt một phía), \texttt{entitlement} (lương, BHXH, giờ làm), \texttt{procedure} (thử việc, giải quyết tranh chấp)
        \item Truy vết multi-hop trên entity L1 (\texttt{Dieu}, \texttt{VanBan}) và L2 (\texttt{ChuThe}, \texttt{TienLuong}, \texttt{CheTai}, \ldots) qua quan hệ \texttt{entitles}, \texttt{obligates}, \texttt{penalizes}, \texttt{guided\_by}, \ldots
        \item Mở rộng láng giềng entity thêm 2 hop trên parquet (\texttt{RELATED\_TO} trong file quan hệ)
        \item Đầu ra: \texttt{MappedLawSnippet.rag\_answer} --- chuỗi suy luận + danh sách Điều trích dẫn (\texttt{cited\_articles})
    \end{itemize}
    \item \textbf{Mapper dự phòng} khi optimized mapper lỗi:
    \begin{itemize}[nosep, leftmargin=*]
        \item Bước~1: \texttt{basic\_search} --- vector trên \texttt{text\_units} (Naive RAG, không traversal đồ thị)
        \item Bước~2 (tuỳ chọn): seed \texttt{Entity.id} theo overlap token với khoản HĐ (\texttt{entity\_seed.py}) $\rightarrow$ Cypher \texttt{RELATED\_TO*1..hops} trên Neo4j (\texttt{neo4j\_kg\_expand.py})
    \end{itemize}
    \item \textbf{Thông tin khai thác từ đồ thị} (khác với chat Local/Global Search):
    \begin{itemize}[nosep, leftmargin=*]
        \item Entity và mô tả ngữ nghĩa (L2): quyền, nghĩa vụ, chế tài, tiền lương, \ldots
        \item Quan hệ có hướng giữa entity --- phục vụ suy luận đa bước (ví dụ: hành vi $\rightarrow$ \texttt{obligates} $\rightarrow$ quy định $\rightarrow$ \texttt{penalizes} $\rightarrow$ chế tài)
        \item Điều luật liên quan qua chuỗi reasoning và neighborhood expansion
        \item \emph{Không} dùng community reports (Global Search) trong pipeline HĐLĐ hiện tại
    \end{itemize}
    \item \textbf{Luồng dữ liệu tóm tắt}: khoản HĐ $\rightarrow$ (rule VR) $\rightarrow$ [nếu phức tạp] map đồ thị $\rightarrow$ ngữ cảnh pháp lý $\rightarrow$ LLM review $\rightarrow$ báo cáo
    \item Tham số pipeline: \texttt{use\_optimized\_mapper=True}, \texttt{skip\_graph\_mapping}, \texttt{use\_neo4j\_knowledge\_graph}, \texttt{neo4j\_graph\_hops} (1--3)
\end{itemize}

\subsection{Quy tắc VR001--VR016}
\label{subsection:3.5.3}
\begin{itemize}[nosep, leftmargin=*]
    \item VR001: Lương $<$ mức tối thiểu vùng (BLLĐ Điều 90 + NĐ 74/2024)
    \item VR002--VR007: Thử việc, giờ làm, phạt tiền NLĐ, thời hạn HĐLĐ, \ldots
    \item VR008--VR016: BHXH, nghỉ phép, BHTN, điều kiện làm việc, \ldots
    \item \texttt{RuleContext}: vùng lương (I--IV), \texttt{max\_probation\_days}
\end{itemize}

\section{Thiết kế giao diện và API}
\label{section:3.6}

\subsection{API FastAPI}
\label{subsection:3.6.1}
\begin{itemize}[nosep, leftmargin=*]
    \item Routes: \texttt{/api/chat}, \texttt{/api/documents}, \texttt{/api/health}
    \item \texttt{ChatRequest}: \texttt{question}, \texttt{mode}, \texttt{domain}
    \item \texttt{ChatResponse}: \texttt{answer}, \texttt{article\_citations}, \texttt{entities\_used}, \texttt{temporal\_warnings}
    \item \texttt{api/services/graph\_loader.py} --- load parquet artifacts
\end{itemize}

\subsection{Frontend Next.js}
\label{subsection:3.6.2}
\begin{itemize}[nosep, leftmargin=*]
    \item UI chat hỏi đáp, quản lý tài liệu
    \item Luồng: nhập câu hỏi $\rightarrow$ chọn mode $\rightarrow$ nhận câu trả lời + trích dẫn
\end{itemize}

\end{document}
